\section{S}

% sagen (to say)
\begin{verbentry}{
  style = {acc},
  toc = {sagen (to say)},
  name = {sagen},
  english = {to say},
  type = {regular, + Akkusativ},
  auxiliary = {haben}
}
 
    
     \vspace{5pt}

    \hrule
    
    \vspace{5pt}

  \begin{tabular}{rl}
     \multicolumn{2}{l}{\textbf{PRÄSENS} }\\
    ich & \textbf{sage}\\
    du & \textbf{sagst}\\
    er/sie/es & \textbf{sagt}\\
    wir & \textbf{sagen}\\
    ihr & \textbf{sagt}\\
    sie/Sie & \textbf{sagen}\\
  \end{tabular}
  \hfill
     \begin{tabular}{rl}
    \multicolumn{2}{l}{\textbf{PRÄTERITUM} }\\
    ich & \textbf{sagte}\\
    du & \textbf{sagtest}\\
    er/sie/es & \textbf{sagte}\\
    wir & \textbf{sagten}\\
    ihr & \textbf{sagtet}\\
    sie/Sie & \textbf{sagten}\\
  \end{tabular}
  \hfill
  \begin{tabular}{rl}
     \multicolumn{2}{l}{\textbf{IMPERATIV} }\\
   & \textbf{sag(e) (du)!}\\
   & \textbf{sagen wir!}\\
   & \textbf{sagt (ihr)!}\\
   & \textbf{sagen Sie!}\\
   & \vspace{10pt}
  \end{tabular}

    \vspace{10pt}

 \begin{tabular}{rl}
     \multicolumn{2}{l}{\textbf{PERFEKT}}\\
   & haben + \textbf{gesagt}\\
  \end{tabular}
\hfill
   \begin{tabular}{rl}
   
    \multicolumn{2}{l}{\textbf{FUTURE I} }\\
  &  werden + \textbf{sagen}\\
 
  \end{tabular}
\hfill
   \begin{tabular}{rl}
      \multicolumn{2}{l}{\textbf{PLUSQUAMPERFEKT}}\\
  &  hatten + \textbf{gesagt}\\
  \end{tabular}
  
  
\vspace{10pt}

\begin{tabular}{lll}
\multicolumn{3}{l}{\textbf{MIT PRÄPOSITIONEN}}\\
& \textbf{---} & \\
\end{tabular}
\hfill
\begin{tabular}{lll}
\multicolumn{3}{l}{\textbf{TRENNBARE VERBEN}}\\
& \textbf{---} & \\
\end{tabular}
\vspace{-5pt}
\end{verbentry}

% sammeln (to collect / gather)
\begin{verbentry}{
  style = {acc},
  toc = {sammeln (to collect / gather)},
  name = {sammeln},
  english = {to collect / gather},
  type = {regular},
  auxiliary = {haben}
}
 
    
     \vspace{5pt}

    \hrule
    
    \vspace{5pt}

  \begin{tabular}{rl}
     \multicolumn{2}{l}{\textbf{PRÄSENS} }\\
    ich & \textbf{sammle}\\
    du & \textbf{sammelst}\\
    er/sie/es & \textbf{sammelt}\\
    wir & \textbf{sammeln}\\
    ihr & \textbf{sammelt}\\
    sie/Sie & \textbf{sammeln}\\
  \end{tabular}
  \hfill
  \begin{tabular}{rl}
    \multicolumn{2}{l}{\textbf{PRÄTERITUM} }\\
    ich & \textbf{sammelte}\\
    du & \textbf{sammeltest}\\
    er/sie/es & \textbf{sammelte}\\
    wir & \textbf{sammelten}\\
    ihr & \textbf{sammeltet}\\
    sie/Sie & \textbf{sammelten}\\
  \end{tabular}
  \hfill
  \begin{tabular}{rl}
     \multicolumn{2}{l}{\textbf{IMPERATIV} }\\
   & \textbf{sammle (du)!}\\
   & \textbf{sammeln wir!}\\
   & \textbf{sammelt (ihr)!}\\
   & \textbf{sammeln Sie!}\\
   & \vspace{10pt}
  \end{tabular}

 \vspace{10pt}

 \begin{tabular}{rl}
     \multicolumn{2}{l}{\textbf{PERFEKT}}\\
   & haben + \textbf{gesammelt}\\
  \end{tabular}
\hfill
\begin{tabular}{rl}
    \multicolumn{2}{l}{\textbf{FUTURE I}}\\
   & werden + \textbf{sammeln}\\
  \end{tabular}
\hfill
\begin{tabular}{rl}
    \multicolumn{2}{l}{\textbf{PLUSQUAMPERFEKT}}\\
   & hatten + \textbf{gesammelt}\\
  \end{tabular}

 \vspace{10pt}

\begin{tabular}{lll}
      \multicolumn{3}{l}{\textbf{MIT PRÄPOSITIONEN}}\\
& \textbf{sammeln für} + Akk & (to collect for / gather for)\\
& \textbf{sammeln von} + Dat & (to collect from / gather from)\\
& \textbf{sammeln in} + Dat & (to collect in / within)\\
\end{tabular}
  \hfill
\begin{tabular}{lll}
      \multicolumn{3}{l}{\textbf{TRENNBARE VERBEN}}\\
 & \textbf{zusammen$|$sammeln} & (to gather together / collect together)\\
 & \textbf{ein$|$sammeln} & (to collect in / pick up)\\
\vspace{10pt}
  \end{tabular}

  \vspace{-5pt}
\end{verbentry}

% saugen (to vacuum / suck)
\begin{verbentry}{
  style = {acc},
  toc = {saugen (to vacuum / suck)},
  name = {saugen},
  english = {to vacuum / suck},
  type = {regular, + Akkusativ},
  auxiliary = {haben}
}


\vspace{5pt}
\hrule
\vspace{5pt}

\begin{tabular}{rl}
\multicolumn{2}{l}{\textbf{PRÄSENS}}\\
ich & \textbf{sauge}\\
du & \textbf{saugst}\\
er/sie/es & \textbf{saugt}\\
wir & \textbf{saugen}\\
ihr & \textbf{saugt}\\
sie/Sie & \textbf{saugen}\\
\end{tabular}
\hfill
\begin{tabular}{rl}
\multicolumn{2}{l}{\textbf{PRÄTERITUM}}\\
ich & \textbf{saugte}\\
du & \textbf{saugtest}\\
er/sie/es & \textbf{saugte}\\
wir & \textbf{saugten}\\
ihr & \textbf{saugtet}\\
sie/Sie & \textbf{saugten}\\
\end{tabular}
\hfill
\begin{tabular}{rl}
\multicolumn{2}{l}{\textbf{IMPERATIV}}\\
& \textbf{saug(e) (du)!}\\
& \textbf{saugen wir!}\\
& \textbf{saugt (ihr)!}\\
& \textbf{saugen Sie!}\\
\end{tabular}

\vspace{10pt}

\begin{tabular}{rl}
\multicolumn{2}{l}{\textbf{PERFEKT}}\\
& haben + \textbf{gesaugt}\\
\end{tabular}
\hfill
\begin{tabular}{rl}
\multicolumn{2}{l}{\textbf{FUTURE I}}\\
& werden + \textbf{saugen}\\
\end{tabular}
\hfill
\begin{tabular}{rl}
\multicolumn{2}{l}{\textbf{PLUSQUAMPERFEKT}}\\
& hatten + \textbf{gesaugt}\\
\end{tabular}

\vspace{10pt}

\begin{tabular}{lll}
\multicolumn{3}{l}{\textbf{MIT PRÄPOSITIONEN}}\\
& \textbf{---} & \\
\end{tabular}
\hfill
\begin{tabular}{lll}
\multicolumn{3}{l}{\textbf{TRENNBARE VERBEN}}\\
& \textbf{auf$|$saugen} & (to absorb)\\
& \textbf{ein$|$saugen} & (to inhale / absorb)\\
\end{tabular}

\vspace{-5pt}
\end{verbentry}

% schaden (to harm / damage)
\begin{verbentry}{
  style = {default},
  toc = {schaden (to harm / damage)},
  name = {schaden},
  english = {to harm / damage},
  type = {regular},
  auxiliary = {haben}
}


\vspace{5pt}
\hrule
\vspace{5pt}

\begin{tabular}{rl}
\multicolumn{2}{l}{\textbf{PRÄSENS}}\\
ich & \textbf{schade}\\
du & \textbf{schadest}\\
er/sie/es & \textbf{schadet}\\
wir & \textbf{schaden}\\
ihr & \textbf{schadet}\\
sie/Sie & \textbf{schaden}\\
\end{tabular}
\hfill
\begin{tabular}{rl}
\multicolumn{2}{l}{\textbf{PRÄTERITUM}}\\
ich & \textbf{schadete}\\
du & \textbf{schadetest}\\
er/sie/es & \textbf{schadete}\\
wir & \textbf{schadeten}\\
ihr & \textbf{schadetet}\\
sie/Sie & \textbf{schadeten}\\
\end{tabular}
\hfill
\begin{tabular}{rl}
\multicolumn{2}{l}{\textbf{IMPERATIV}}\\
& \textbf{schade (du)!}\\
& \textbf{schaden wir!}\\
& \textbf{schadet (ihr)!}\\
& \textbf{schaden Sie!}\\
\end{tabular}

\vspace{10pt}

\begin{tabular}{rl}
\multicolumn{2}{l}{\textbf{PERFEKT}}\\
& haben + \textbf{geschadet}\\
\end{tabular}
\hfill
\begin{tabular}{rl}
\multicolumn{2}{l}{\textbf{FUTURE I}}\\
& werden + \textbf{schaden}\\
\end{tabular}
\hfill
\begin{tabular}{rl}
\multicolumn{2}{l}{\textbf{PLUSQUAMPERFEKT}}\\
& hatten + \textbf{geschadet}\\
\end{tabular}

\vspace{10pt}

\begin{tabular}{lll}
\multicolumn{3}{l}{\textbf{MIT PRÄPOSITIONEN}}\\
& \textbf{schaden durch} + Akk & (to be harmed by)\\
\end{tabular}
\hfill
\begin{tabular}{lll}
\multicolumn{3}{l}{\textbf{TRENNBARE VERBEN}}\\
& \textbf{---} & \\
\end{tabular}
\vspace{-5pt}
\end{verbentry}

% schaffen (to manage, create)
\begin{verbentry}{
  style = {acc},
  toc = {schaffen (to manage, create)},
  name = {schaffen},
  english = {to manage, create},
  type = {irregular, + Akkusativ},
  auxiliary = {haben}
}
 
    
     \vspace{5pt}

    \hrule
    
    \vspace{5pt}

  \begin{tabular}{rl}
     \multicolumn{2}{l}{\textbf{PRÄSENS} }\\
    ich & \textbf{schaffe}\\
    du & \textbf{schaffst}\\
    er/sie/es & \textbf{schafft}\\
    wir & \textbf{schaffen}\\
    ihr & \textbf{schafft}\\
    sie/Sie & \textbf{schaffen}\\
  \end{tabular}
  \hfill
     \begin{tabular}{rl}
    \multicolumn{2}{l}{\textbf{PRÄTERITUM} }\\
    ich & \textbf{schuf}\\
    du & \textbf{schufst}\\
    er/sie/es & \textbf{schuf}\\
    wir & \textbf{schufen}\\
    ihr & \textbf{schuft}\\
    sie/Sie & \textbf{schufen}\\
  \end{tabular}
  \hfill
  \begin{tabular}{rl}
     \multicolumn{2}{l}{\textbf{IMPERATIV} }\\
   & \textbf{schaff(e) (du)!}\\
   & \textbf{schaffen wir!}\\
   & \textbf{schafft (ihr)!}\\
   & \textbf{schaffen Sie!}\\
   & \vspace{10pt}
  \end{tabular}

    \vspace{10pt}

 \begin{tabular}{rl}
     \multicolumn{2}{l}{\textbf{PERFEKT}}\\
   & haben + \textbf{geschaffen}\\
  \end{tabular}
\hfill
   \begin{tabular}{rl}
   
    \multicolumn{2}{l}{\textbf{FUTURE I} }\\
  &  werden + \textbf{schaffen}\\
 
  \end{tabular}
\hfill
   \begin{tabular}{rl}
      \multicolumn{2}{l}{\textbf{PLUSQUAMPERFEKT}}\\
  &  hatten + \textbf{geschaffen}\\
  \end{tabular}
  
  
\vspace{10pt}

\begin{tabular}{lll}
\multicolumn{3}{l}{\textbf{MIT PRÄPOSITIONEN}}\\
& \textbf{---} & \\
\end{tabular}
\hfill
\begin{tabular}{lll}
\multicolumn{3}{l}{\textbf{TRENNBARE VERBEN}}\\
& \textbf{---} & \\
\end{tabular}
\vspace{-5pt}
\end{verbentry}

% schalten (to switch)
\begin{verbentry}{
  style = {default},
  toc = {schalten (to switch)},
  name = {schalten},
  english = {to switch},
  type = {regular},
  auxiliary = {haben}
}
 

\vspace{5pt}
\hrule
\vspace{5pt}

\begin{tabular}{rl}
\multicolumn{2}{l}{\textbf{PRÄSENS}}\\
ich & \textbf{schalte}\\
du & \textbf{schaltest}\\
er/sie/es & \textbf{schaltet}\\
wir & \textbf{schalten}\\
ihr & \textbf{schaltet}\\
sie/Sie & \textbf{schalten}\\
\end{tabular}
\hfill
\begin{tabular}{rl}
\multicolumn{2}{l}{\textbf{PRÄTERITUM}}\\
ich & \textbf{schaltete}\\
du & \textbf{schaltetest}\\
er/sie/es & \textbf{schaltete}\\
wir & \textbf{schalteten}\\
ihr & \textbf{schaltetet}\\
sie/Sie & \textbf{schalteten}\\
\end{tabular}
\hfill
\begin{tabular}{rl}
\multicolumn{2}{l}{\textbf{IMPERATIV}}\\
& \textbf{schalte (du)!}\\
& \textbf{schalten wir!}\\
& \textbf{schaltet (ihr)!}\\
& \textbf{schalten Sie!}\\
\end{tabular}

\vspace{10pt}

\begin{tabular}{rl}
\multicolumn{2}{l}{\textbf{PERFEKT}}\\
& haben + \textbf{geschaltet}\\
\end{tabular}
\hfill
\begin{tabular}{rl}
\multicolumn{2}{l}{\textbf{FUTURE I}}\\
& werden + \textbf{schalten}\\
\end{tabular}
\hfill
\begin{tabular}{rl}
\multicolumn{2}{l}{\textbf{PLUSQUAMPERFEKT}}\\
& hatten + \textbf{geschaltet}\\
\end{tabular}

\vspace{10pt}

\begin{tabular}{lll}
\multicolumn{3}{l}{\textbf{MIT PRÄPOSITIONEN}}\\
& \textbf{schalten auf} + Akk & (to switch to)\\
\vspace{10pt}
\end{tabular}
\hfill
\begin{tabular}{lll}
\multicolumn{3}{l}{\textbf{TRENNBARE VERBEN}}\\
& \textbf{ein$|$schalten} & (to switch on)\\
& \textbf{aus$|$schalten} & (to switch off)\\
\end{tabular}

\vspace{-5pt}
\end{verbentry}

% schauen (to look / watch)
\begin{verbentry}{
  style = {default},
  toc = {schauen (to look / watch)},
  name = {schauen},
  english = {to look / watch},
  type = {regular},
  auxiliary = {haben}
}
 

\vspace{5pt}
\hrule
\vspace{5pt}

\begin{tabular}{rl}
\multicolumn{2}{l}{\textbf{PRÄSENS}}\\
ich & \textbf{schaue}\\
du & \textbf{schaust}\\
er/sie/es & \textbf{schaut}\\
wir & \textbf{schauen}\\
ihr & \textbf{schaut}\\
sie/Sie & \textbf{schauen}\\
\end{tabular}
\hfill
\begin{tabular}{rl}
\multicolumn{2}{l}{\textbf{PRÄTERITUM}}\\
ich & \textbf{schaute}\\
du & \textbf{schautest}\\
er/sie/es & \textbf{schaute}\\
wir & \textbf{schauten}\\
ihr & \textbf{schautet}\\
sie/Sie & \textbf{schauten}\\
\end{tabular}
\hfill
\begin{tabular}{rl}
\multicolumn{2}{l}{\textbf{IMPERATIV}}\\
& \textbf{schaue (du)!}\\
& \textbf{schauen wir!}\\
& \textbf{schaut (ihr)!}\\
& \textbf{schauen Sie!}\\
\end{tabular}

\vspace{10pt}

\begin{tabular}{rl}
\multicolumn{2}{l}{\textbf{PERFEKT}}\\
& haben + \textbf{geschaut}\\
\end{tabular}
\hfill
\begin{tabular}{rl}
\multicolumn{2}{l}{\textbf{FUTURE I}}\\
& werden + \textbf{schauen}\\
\end{tabular}
\hfill
\begin{tabular}{rl}
\multicolumn{2}{l}{\textbf{PLUSQUAMPERFEKT}}\\
& hatten + \textbf{geschaut}\\
\end{tabular}

\vspace{10pt}

\begin{tabular}{lll}
\multicolumn{3}{l}{\textbf{MIT PRÄPOSITIONEN}}\\
& \textbf{schauen auf} + Akk & (to look at)\\
\end{tabular}
\hfill
\begin{tabular}{lll}
\multicolumn{3}{l}{\textbf{TRENNBARE VERBEN}}\\
& \textbf{an$|$schauen} & (to look at)\\
\end{tabular}

\vspace{-5pt}
\end{verbentry}

% scheinen (to shine)
\begin{verbentry}{
  style = {default},
  toc = {scheinen (to shine)},
  name = {scheinen},
  english = {to shine},
  type = {irregular},
  auxiliary = {haben}
}
 
    
     \vspace{5pt}

    \hrule
    
    \vspace{5pt}

  \begin{tabular}{rl}
     \multicolumn{2}{l}{\textbf{PRÄSENS} }\\
    ich & \textbf{scheine}\\
    du & \textbf{scheinst}\\
    er/sie/es & \textbf{scheint}\\
    wir & \textbf{scheinen}\\
    ihr & \textbf{scheint}\\
    sie/Sie & \textbf{scheinen}\\
  \end{tabular}
  \hfill
     \begin{tabular}{rl}
    \multicolumn{2}{l}{\textbf{PRÄTERITUM} }\\
    ich & \textbf{schien}\\
    du & \textbf{schienst}\\
    er/sie/es & \textbf{schien}\\
    wir & \textbf{schienen}\\
    ihr & \textbf{schient}\\
    sie/Sie & \textbf{schienen}\\
  \end{tabular}
  \hfill
  \begin{tabular}{rl}
     \multicolumn{2}{l}{\textbf{IMPERATIV} }\\
   & \textbf{schein(e) (du)!}\\
   & \textbf{scheinen wir!}\\
   & \textbf{scheint (ihr)!}\\
   & \textbf{scheinen Sie!}\\
   & \vspace{10pt}
  \end{tabular}

    \vspace{10pt}

 \begin{tabular}{rl}
     \multicolumn{2}{l}{\textbf{PERFEKT}}\\
   & haben + \textbf{geschienen}\\
  \end{tabular}
\hfill
   \begin{tabular}{rl}
   
    \multicolumn{2}{l}{\textbf{FUTURE I} }\\
  &  werden + \textbf{scheinen}\\
 
  \end{tabular}
\hfill
   \begin{tabular}{rl}
      \multicolumn{2}{l}{\textbf{PLUSQUAMPERFEKT}}\\
  &  hatten + \textbf{geschienen}\\
  \end{tabular}
  
  \vspace{10pt}



\begin{tabular}{lll}
\multicolumn{3}{l}{\textbf{MIT PRÄPOSITIONEN}}\\
& \textbf{---} & \\
\end{tabular}
\hfill
\begin{tabular}{lll}
\multicolumn{3}{l}{\textbf{TRENNBARE VERBEN}}\\
& \textbf{---} & \\
\end{tabular}

\vspace{-5pt}
\end{verbentry}

% schicken (to send)
\begin{verbentry}{
  style = {default},
  toc = {schicken (to send)},
  name = {schicken},
  english = {to send},
  type = {regular},
  auxiliary = {haben}
}


\vspace{5pt}
\hrule
\vspace{5pt}

\begin{tabular}{rl}
\multicolumn{2}{l}{\textbf{PRÄSENS}}\\
ich & \textbf{schicke}\\
du & \textbf{schickst}\\
er/sie/es & \textbf{schickt}\\
wir & \textbf{schicken}\\
ihr & \textbf{schickt}\\
sie/Sie & \textbf{schicken}\\
\end{tabular}
\hfill
\begin{tabular}{rl}
\multicolumn{2}{l}{\textbf{PRÄTERITUM}}\\
ich & \textbf{schickte}\\
du & \textbf{schicktest}\\
er/sie/es & \textbf{schickte}\\
wir & \textbf{schickten}\\
ihr & \textbf{schicktet}\\
sie/Sie & \textbf{schickten}\\
\end{tabular}
\hfill
\begin{tabular}{rl}
\multicolumn{2}{l}{\textbf{IMPERATIV}}\\
& \textbf{schick(e) (du)!}\\
& \textbf{schicken wir!}\\
& \textbf{schickt (ihr)!}\\
& \textbf{schicken Sie!}\\
\end{tabular}

\vspace{10pt}

\begin{tabular}{rl}
\multicolumn{2}{l}{\textbf{PERFEKT}}\\
& haben + \textbf{geschickt}\\
\end{tabular}
\hfill
\begin{tabular}{rl}
\multicolumn{2}{l}{\textbf{FUTURE I}}\\
& werden + \textbf{schicken}\\
\end{tabular}
\hfill
\begin{tabular}{rl}
\multicolumn{2}{l}{\textbf{PLUSQUAMPERFEKT}}\\
& hatten + \textbf{geschickt}\\
\end{tabular}

\vspace{10pt}

\begin{tabular}{lll}
\multicolumn{3}{l}{\textbf{MIT PRÄPOSITIONEN}}\\
& \textbf{schicken an} + Akk & (to send to)\\
& \textbf{schicken mit} + Dat & (to send with)\\
\end{tabular}
\hfill
\begin{tabular}{lll}
\multicolumn{3}{l}{\textbf{TRENNBARE VERBEN}}\\
& \textbf{ab$|$schicken} & (to send off)\\
& \textbf{ein$|$schicken} & (to submit / send in)\\
& \textbf{zurück$|$schicken} & (to send back)\\
\end{tabular}
\vspace{-5pt}
\end{verbentry}

% schlafen (to sleep)
\begin{verbentry}{
  style = {default},
  toc = {schlafen (to sleep)},
  name = {schlafen},
  english = {to sleep},
  type = {irregular},
  auxiliary = {haben}
}
 
    
     \vspace{5pt}

    \hrule
    
    \vspace{5pt}

  \begin{tabular}{rl}
     \multicolumn{2}{l}{\textbf{PRÄSENS} }\\
    ich & \textbf{schlafe}\\
    du & \textbf{schläfst}\\
    er/sie/es & \textbf{schläft}\\
    wir & \textbf{schlafen}\\
    ihr & \textbf{schlaft}\\
    sie/Sie & \textbf{schlafen}\\
  \end{tabular}
  \hfill
     \begin{tabular}{rl}
    \multicolumn{2}{l}{\textbf{PRÄTERITUM} }\\
    ich & \textbf{schlief}\\
    du & \textbf{schliefest}\\
    er/sie/es & \textbf{schlief}\\
    wir & \textbf{schliefen}\\
    ihr & \textbf{schlieft}\\
    sie/Sie & \textbf{schlieften}\\
  \end{tabular}
  \hfill
  \begin{tabular}{rl}
     \multicolumn{2}{l}{\textbf{IMPERATIV} }\\
   & \textbf{schlaf(e) (du)!}\\
   & \textbf{schlafen wir!}\\
   & \textbf{schlaft (ihr)!}\\
   & \textbf{schlafen Sie!}\\
   & \vspace{10pt}
  \end{tabular}

    \vspace{10pt}

 \begin{tabular}{rl}
     \multicolumn{2}{l}{\textbf{PERFEKT}}\\
  &  haben + \textbf{geschlafen}\\
  \end{tabular}
\hfill
   \begin{tabular}{rl}
   
    \multicolumn{2}{l}{\textbf{FUTURE I} }\\
  &  werden + \textbf{schlafen}\\
 
  \end{tabular}
\hfill
   \begin{tabular}{rl}
      \multicolumn{2}{l}{\textbf{PLUSQUAMPERFEKT}}\\
  &  hatten + \textbf{geschlafen}\\
  \end{tabular}
  
  \vspace{10pt}



\begin{tabular}{lll}
\multicolumn{3}{l}{\textbf{MIT PRÄPOSITIONEN}}\\
& \textbf{---} & \\
\end{tabular}
\hfill
\begin{tabular}{lll}
\multicolumn{3}{l}{\textbf{TRENNBARE VERBEN}}\\
& \textbf{---} & \\
\end{tabular}

  \vspace{-5pt}
\end{verbentry}

% schlagen (to hit / beat / strike)
\begin{verbentry}{
  style = {acc},
  toc = {schlagen (to hit / beat / strike)},
  name = {schlagen},
  english = {to hit / beat / strike},
  type = {irregular, + Akkusativ},
  auxiliary = {haben}
}
 
    
     \vspace{5pt}

    \hrule
    
    \vspace{5pt}

  \begin{tabular}{rl}
     \multicolumn{2}{l}{\textbf{PRÄSENS} }\\
    ich & \textbf{schlage}\\
    du & \textbf{schlägst}\\
    er/sie/es & \textbf{schlägt}\\
    wir & \textbf{schlagen}\\
    ihr & \textbf{schlagt}\\
    sie/Sie & \textbf{schlagen}\\
  \end{tabular}
  \hfill
  \begin{tabular}{rl}
    \multicolumn{2}{l}{\textbf{PRÄTERITUM} }\\
    ich & \textbf{schlug}\\
    du & \textbf{schlugst}\\
    er/sie/es & \textbf{schlug}\\
    wir & \textbf{schlugen}\\
    ihr & \textbf{schlugt}\\
    sie/Sie & \textbf{schlugen}\\
  \end{tabular}
  \hfill
  \begin{tabular}{rl}
     \multicolumn{2}{l}{\textbf{IMPERATIV} }\\
   & \textbf{schlag (du)!}\\
   & \textbf{schlagen wir!}\\
   & \textbf{schlagt (ihr)!}\\
   & \textbf{schlagen Sie!}\\
   & \vspace{10pt}
  \end{tabular}

 \vspace{10pt}

 \begin{tabular}{rl}
     \multicolumn{2}{l}{\textbf{PERFEKT}}\\
   & haben + \textbf{geschlagen}\\
  \end{tabular}
\hfill
\begin{tabular}{rl}
    \multicolumn{2}{l}{\textbf{FUTURE I}}\\
   & werden + \textbf{schlagen}\\
  \end{tabular}
\hfill
\begin{tabular}{rl}
    \multicolumn{2}{l}{\textbf{PLUSQUAMPERFEKT}}\\
   & hatten + \textbf{geschlagen}\\
  \end{tabular}

 \vspace{10pt}

\begin{tabular}{lll}
      \multicolumn{3}{l}{\textbf{MIT PRÄPOSITIONEN}}\\
& \textbf{schlagen auf} + Akk & (to hit on / strike)\\
& \textbf{schlagen gegen} + Akk & (to hit against / collide)\\
& \textbf{schlagen für} + Akk & (to vote / advocate for)\\
\vspace{20pt}
\end{tabular}
  \hfill
\begin{tabular}{lll}
      \multicolumn{3}{l}{\textbf{TRENNBARE VERBEN}}\\
 & \textbf{auf$|$schlagen} & (to hit open / pitch / unfold)\\
 & \textbf{ein$|$schlagen} & (to break / strike / hammer in)\\
 & \textbf{vor$|$schlagen} & (to propose / suggest)\\
 & \textbf{zurück$|$schlagen} & (to hit back / retaliate)\\
 & \textbf{zusammen$|$schlagen} & (to beat / knock together)\\
  \end{tabular}

  \vspace{-5pt}
\end{verbentry}

% schleifen (to sand / grind / sharpen)
\begin{verbentry}{
  style = {default},
  toc = {schleifen (to sand / grind / sharpen)},
  name = {schleifen},
  english = {to sand / grind / sharpen},
  type = {regular},
  auxiliary = {haben}
}
 
    
     \vspace{5pt}

    \hrule
    
    \vspace{5pt}

  \begin{tabular}{rl}
     \multicolumn{2}{l}{\textbf{PRÄSENS} }\\
    ich & \textbf{schleife}\\
    du & \textbf{schleifst}\\
    er/sie/es & \textbf{schleift}\\
    wir & \textbf{schleifen}\\
    ihr & \textbf{schleift}\\
    sie/Sie & \textbf{schleifen}\\
  \end{tabular}
  \hfill
  \begin{tabular}{rl}
    \multicolumn{2}{l}{\textbf{PRÄTERITUM} }\\
    ich & \textbf{schliff}\\
    du & \textbf{schliffst}\\
    er/sie/es & \textbf{schliff}\\
    wir & \textbf{schliffen}\\
    ihr & \textbf{schlifft}\\
    sie/Sie & \textbf{schliffen}\\
  \end{tabular}
  \hfill
  \begin{tabular}{rl}
     \multicolumn{2}{l}{\textbf{IMPERATIV} }\\
   & \textbf{schleife (du)!}\\
   & \textbf{schleifen wir!}\\
   & \textbf{schleift (ihr)!}\\
   & \textbf{schleifen Sie!}\\
   & \vspace{10pt}
  \end{tabular}

 \vspace{10pt}

 \begin{tabular}{rl}
     \multicolumn{2}{l}{\textbf{PERFEKT}}\\
   & haben + \textbf{geschliffen}\\
  \end{tabular}
\hfill
\begin{tabular}{rl}
    \multicolumn{2}{l}{\textbf{FUTURE I}}\\
   & werden + \textbf{schleifen}\\
  \end{tabular}
\hfill
\begin{tabular}{rl}
    \multicolumn{2}{l}{\textbf{PLUSQUAMPERFEKT}}\\
   & hatten + \textbf{geschliffen}\\
  \end{tabular}

 \vspace{10pt}

\begin{tabular}{lll}
      \multicolumn{3}{l}{\textbf{MIT PRÄPOSITIONEN}}\\
& \textbf{schleifen an} + Dat & (to sharpen on / work on)\\
& \textbf{schleifen mit} + Dat & (to sand / grind with)\\
\end{tabular}
  \hfill
\begin{tabular}{lll}
      \multicolumn{3}{l}{\textbf{TRENNBARE VERBEN}}\\
 & \textbf{ab$|$schleifen} & (to grind off / remove)\\
 & \textbf{ein$|$schleifen} & (to polish / fit in)\\
  \end{tabular}


\vspace{-5pt}
\end{verbentry}

% schlendern (to stroll)
\begin{verbentry}{
  style = {default},
  toc = {schlendern (to stroll)},
  name = {schlendern},
  english = {to stroll},
  type = {regular},
  auxiliary = {sein}
}


\vspace{5pt}
\hrule
\vspace{5pt}

\begin{tabular}{rl}
\multicolumn{2}{l}{\textbf{PRÄSENS}}\\
ich & \textbf{schlendere}\\
du & \textbf{schlenderst}\\
er/sie/es & \textbf{schlendert}\\
wir & \textbf{schlendern}\\
ihr & \textbf{schlendert}\\
sie/Sie & \textbf{schlendern}\\
\end{tabular}
\hfill
\begin{tabular}{rl}
\multicolumn{2}{l}{\textbf{PRÄTERITUM}}\\
ich & \textbf{schlenderte}\\
du & \textbf{schlendertest}\\
er/sie/es & \textbf{schlenderte}\\
wir & \textbf{schlenderten}\\
ihr & \textbf{schlendertet}\\
sie/Sie & \textbf{schlenderten}\\
\end{tabular}
\hfill
\begin{tabular}{rl}
\multicolumn{2}{l}{\textbf{IMPERATIV}}\\
& \textbf{schlendere (du)!}\\
& \textbf{schlendern wir!}\\
& \textbf{schlendert (ihr)!}\\
& \textbf{schlendern Sie!}\\
\end{tabular}

\vspace{10pt}

\begin{tabular}{rl}
\multicolumn{2}{l}{\textbf{PERFEKT}}\\
& sein + \textbf{geschlendert}\\
\end{tabular}
\hfill
\begin{tabular}{rl}
\multicolumn{2}{l}{\textbf{FUTURE I}}\\
& werden + \textbf{schlendern}\\
\end{tabular}
\hfill
\begin{tabular}{rl}
\multicolumn{2}{l}{\textbf{PLUSQUAMPERFEKT}}\\
& waren + \textbf{geschlendert}\\
\end{tabular}

\vspace{10pt}

\begin{tabular}{lll}
\multicolumn{3}{l}{\textbf{MIT PRÄPOSITIONEN}}\\
& \textbf{schlendern durch} + Akk & (to stroll through)\\
& \textbf{schlendern über} + Akk & (to stroll across)\\
\end{tabular}
\hfill
\begin{tabular}{lll}
\multicolumn{3}{l}{\textbf{TRENNBARE VERBEN}}\\
& \textbf{---} & \\
\end{tabular}

\vspace{-5pt}
\end{verbentry}

% schließen (to close)
\begin{verbentry}{
  style = {acc},
  toc = {schließen (to close)},
  name = {schließen},
  english = {to close},
  type = {irregular, + Akkusativ},
  auxiliary = {haben}
}
 
    
     \vspace{5pt}

    \hrule
    
    \vspace{5pt}

  \begin{tabular}{rl}
     \multicolumn{2}{l}{\textbf{PRÄSENS} }\\
    ich & \textbf{schließe}\\
    du & \textbf{schließt}\\
    er/sie/es & \textbf{schließt}\\
    wir & \textbf{schließen}\\
    ihr & \textbf{schließt}\\
    sie/Sie & \textbf{schließen}\\
  \end{tabular}
  \hfill
     \begin{tabular}{rl}
    \multicolumn{2}{l}{\textbf{PRÄTERITUM} }\\
    ich & \textbf{schloss}\\
    du & \textbf{schlossest}\\
    er/sie/es & \textbf{schloss}\\
    wir & \textbf{schlossen}\\
    ihr & \textbf{schlosst}\\
    sie/Sie & \textbf{schlossen}\\
  \end{tabular}
  \hfill
  \begin{tabular}{rl}
     \multicolumn{2}{l}{\textbf{IMPERATIV} }\\
   & \textbf{schließ(e) (du)!}\\
   & \textbf{schließen wir!}\\
   & \textbf{schließt (ihr)!}\\
   & \textbf{schließen Sie!}\\
   & \vspace{10pt}
  \end{tabular}

    \vspace{10pt}

 \begin{tabular}{rl}
     \multicolumn{2}{l}{\textbf{PERFEKT}}\\
   & haben + \textbf{geschlossen}\\
  \end{tabular}
\hfill
   \begin{tabular}{rl}
   
    \multicolumn{2}{l}{\textbf{FUTURE I} }\\
   & werden + \textbf{schließen}\\
 
  \end{tabular}
\hfill
   \begin{tabular}{rl}
      \multicolumn{2}{l}{\textbf{PLUSQUAMPERFEKT}}\\
   & hatten + \textbf{geschlossen}\\
  \end{tabular}
  
  \vspace{10pt}

   \begin{tabular}{lll}
      \multicolumn{3}{l}{\textbf{MIT PRÄPOSITIONEN}}\\
      & \textbf{schließen auf} + Akk & (to infer / conclude)\\
&\textbf{schließen aus} + Dat & (to infer from)\\
&\textbf{schließen mit} + Dat & (to close with / end with)\\
& \textbf{anschließen an} + Akk & (to conect to)\\
& \textbf{abschließen mit} + Dat & (to finish with)\\
&\textbf{einschließen in} + Akk & (to include in)\\
&\textbf{ausschließen von} + Dat & (to exlude from)
  \end{tabular}
  \hfill
   \begin{tabular}{lll}
      \multicolumn{3}{l}{\textbf{TRENNBARE VERBEN}}\\
 & \textbf{ab$|$schließen} & (to lock)\\
 & \textbf{auf$|$schließen} & (to unlock)\\
 & \textbf{ein$|$schließen} & (to include / enclose)\\
  & \textbf{aus$|$schließen} & (to exclude)\\
  & \textbf{an$|$schließen} & (to connect / join)\\
  & \textbf{zu$|$schließen} & (to lock (shut))
  \end{tabular}




  \vspace{-5pt}
\end{verbentry}

% schmecken (to taste)
\begin{verbentry}{
  style = {dat},
  toc = {schmecken (to taste)},
  name = {schmecken},
  english = {to taste},
  type = {regular, + Akkusativ/Dativ},
  auxiliary = {haben}
}
 
    
     \vspace{5pt}

    \hrule
    
    \vspace{5pt}

  \begin{tabular}{rl}
     \multicolumn{2}{l}{\textbf{PRÄSENS} }\\
    ich & \textbf{schmecke}\\
    du & \textbf{schmeckst}\\
    er/sie/es & \textbf{schmeckt}\\
    wir & \textbf{schmecken}\\
    ihr & \textbf{schmeckt}\\
    sie/Sie & \textbf{schmecken}\\
  \end{tabular}
  \hfill
     \begin{tabular}{rl}
    \multicolumn{2}{l}{\textbf{PRÄTERITUM} }\\
    ich & \textbf{schmeckte}\\
    du & \textbf{schmecktest}\\
    er/sie/es & \textbf{schmeckte}\\
    wir & \textbf{schmeckten}\\
    ihr & \textbf{schmecktet}\\
    sie/Sie & \textbf{schmeckten}\\
  \end{tabular}
  \hfill
  \begin{tabular}{rl}
     \multicolumn{2}{l}{\textbf{IMPERATIV} }\\
   & \textbf{schmeck(e) (du)!}\\
   & \textbf{schmecken wir!}\\
   & \textbf{schmeckt (ihr)!}\\
   & \textbf{schmecken Sie!}\\
   & \vspace{10pt}
  \end{tabular}

    \vspace{10pt}

 \begin{tabular}{rl}
     \multicolumn{2}{l}{\textbf{PERFEKT}}\\
  &  haben + \textbf{geschmeckt}\\
  \end{tabular}
\hfill
   \begin{tabular}{rl}
   
    \multicolumn{2}{l}{\textbf{FUTURE I} }\\
  &  werden + \textbf{schmecken}\\
 
  \end{tabular}
\hfill
   \begin{tabular}{rl}
      \multicolumn{2}{l}{\textbf{PLUSQUAMPERFEKT}}\\
  &  hatten + \textbf{geschmeckt}\\
  \end{tabular}
  
  
\vspace{10pt}

\begin{tabular}{lll}
\multicolumn{3}{l}{\textbf{MIT PRÄPOSITIONEN}}\\
& \textbf{---} & \\
\end{tabular}
\hfill
\begin{tabular}{lll}
\multicolumn{3}{l}{\textbf{TRENNBARE VERBEN}}\\
& \textbf{---} & \\
\end{tabular}
\vspace{-5pt}
\end{verbentry}

% schmücken (to decorate)
\begin{verbentry}{
  style = {default},
  toc = {schmücken (to decorate)},
  name = {schmücken},
  english = {to decorate},
  type = {regular},
  auxiliary = {haben}
}


\vspace{5pt}
\hrule
\vspace{5pt}

\begin{tabular}{rl}
\multicolumn{2}{l}{\textbf{PRÄSENS}}\\
ich & \textbf{schmücke}\\
du & \textbf{schmückst}\\
er/sie/es & \textbf{schmückt}\\
wir & \textbf{schmücken}\\
ihr & \textbf{schmückt}\\
sie/Sie & \textbf{schmücken}\\
\end{tabular}
\hfill
\begin{tabular}{rl}
\multicolumn{2}{l}{\textbf{PRÄTERITUM}}\\
ich & \textbf{schmückte}\\
du & \textbf{schmücktest}\\
er/sie/es & \textbf{schmückte}\\
wir & \textbf{schmückten}\\
ihr & \textbf{schmücktet}\\
sie/Sie & \textbf{schmückten}\\
\end{tabular}
\hfill
\begin{tabular}{rl}
\multicolumn{2}{l}{\textbf{IMPERATIV}}\\
& \textbf{schmücke (du)!}\\
& \textbf{schmücken wir!}\\
& \textbf{schmückt (ihr)!}\\
& \textbf{schmücken Sie!}\\
\end{tabular}

\vspace{10pt}

\begin{tabular}{rl}
\multicolumn{2}{l}{\textbf{PERFEKT}}\\
& haben + \textbf{geschmückt}\\
\end{tabular}
\hfill
\begin{tabular}{rl}
\multicolumn{2}{l}{\textbf{FUTURE I}}\\
& werden + \textbf{schmücken}\\
\end{tabular}
\hfill
\begin{tabular}{rl}
\multicolumn{2}{l}{\textbf{PLUSQUAMPERFEKT}}\\
& hatten + \textbf{geschmückt}\\
\end{tabular}

\vspace{10pt}

\begin{tabular}{lll}
\multicolumn{3}{l}{\textbf{MIT PRÄPOSITIONEN}}\\
& \textbf{schmücken mit} + Dat & (to decorate with)\\
& \textbf{schmücken für} + Akk & (to decorate for)\\
\end{tabular}
\hfill
\begin{tabular}{lll}
\multicolumn{3}{l}{\textbf{TRENNBARE VERBEN}}\\
& \textbf{aus$|$schmücken} & (to decorate elaborately)\\
\end{tabular}
\vspace{-5pt}
\end{verbentry}

% schneien (to snow)
\begin{verbentry}{
  style = {default},
  toc = {schneien (to snow)},
  name = {schneien},
  english = {to snow},
  type = {irregular},
  auxiliary = {haben}
}
 
    
     \vspace{5pt}

    \hrule
    
    \vspace{5pt}

  \begin{tabular}{rl}
     \multicolumn{2}{l}{\textbf{PRÄSENS} }\\
    ich & \textbf{-}\\
    du & \textbf{-}\\
    es & \textbf{schneit}\\
    wir & \textbf{-}\\
    ihr & \textbf{-}\\
    sie/Sie & \textbf{-}\\
  \end{tabular}
  \hfill
     \begin{tabular}{rl}
    \multicolumn{2}{l}{\textbf{PRÄTERITUM} }\\
    ich & \textbf{-}\\
    du & \textbf{-}\\
    es & \textbf{schneite}\\
    wir & \textbf{-}\\
    ihr & \textbf{-}\\
    sie/Sie & \textbf{-}\\
  \end{tabular}
  \hfill
  \begin{tabular}{rl}
     \multicolumn{2}{l}{\textbf{IMPERATIV} }\\
   & \textbf{-}\\
   & \textbf{-}\\
   & \textbf{-}\\
   & \textbf{-}\\
   & \vspace{10pt}
  \end{tabular}

    \vspace{10pt}

 \begin{tabular}{rl}
     \multicolumn{2}{l}{\textbf{PERFEKT}}\\
  &  haben + \textbf{geschneit}\\
  \end{tabular}
\hfill
   \begin{tabular}{rl}
   
    \multicolumn{2}{l}{\textbf{FUTURE I} }\\
  &  werden + \textbf{schneien}\\
 
  \end{tabular}
\hfill
   \begin{tabular}{rl}
      \multicolumn{2}{l}{\textbf{PLUSQUAMPERFEKT}}\\
  &  hatten + \textbf{geschneit}\\
  \end{tabular}
  
  \vspace{10pt}



\begin{tabular}{lll}
\multicolumn{3}{l}{\textbf{MIT PRÄPOSITIONEN}}\\
& \textbf{---} & \\
\end{tabular}
\hfill
\begin{tabular}{lll}
\multicolumn{3}{l}{\textbf{TRENNBARE VERBEN}}\\
& \textbf{---} & \\
\end{tabular}

  \vspace{-5pt}
\end{verbentry}

% schreiben (to write)
\begin{verbentry}{
  style = {acc},
  toc = {schreiben (to write)},
  name = {schreiben},
  english = {to write},
  type = {irregular, + Akkusativ},
  auxiliary = {haben}
}
 
    
     \vspace{5pt}

    \hrule
    
    \vspace{5pt}

  \begin{tabular}{rl}
     \multicolumn{2}{l}{\textbf{PRÄSENS} }\\
    ich & \textbf{schreibe}\\
    du & \textbf{schreibst}\\
    er/sie/es & \textbf{schreibt}\\
    wir & \textbf{schreiben}\\
    ihr & \textbf{schreibt}\\
    sie/Sie & \textbf{schreiben}\\
  \end{tabular}
  \hfill
     \begin{tabular}{rl}
    \multicolumn{2}{l}{\textbf{PRÄTERITUM} }\\
    ich & \textbf{schrieb}\\
    du & \textbf{schriebst}\\
    er/sie/es & \textbf{schrieb}\\
    wir & \textbf{schrieben}\\
    ihr & \textbf{schriebt}\\
    sie/Sie & \textbf{schrieben}\\
  \end{tabular}
  \hfill
  \begin{tabular}{rl}
     \multicolumn{2}{l}{\textbf{IMPERATIV} }\\
   & \textbf{schreib(e) (du)!}\\
   & \textbf{schreiben wir!}\\
   & \textbf{schreibt (ihr)!}\\
   & \textbf{schreiben Sie!}\\
   & \vspace{10pt}
  \end{tabular}

    \vspace{10pt}

 \begin{tabular}{rl}
     \multicolumn{2}{l}{\textbf{PERFEKT}}\\
  &  haben + \textbf{geschrieben}\\
  \end{tabular}
\hfill
   \begin{tabular}{rl}
   
    \multicolumn{2}{l}{\textbf{FUTURE I} }\\
  &  werden + \textbf{schreiben}\\
 
  \end{tabular}
\hfill
   \begin{tabular}{rl}
      \multicolumn{2}{l}{\textbf{PLUSQUAMPERFEKT}}\\
  &  hatten + \textbf{geschrieben}\\
  \end{tabular}
  
  
\vspace{10pt}

\begin{tabular}{lll}
\multicolumn{3}{l}{\textbf{MIT PRÄPOSITIONEN}}\\
& \textbf{---} & \\
\end{tabular}
\hfill
\begin{tabular}{lll}
\multicolumn{3}{l}{\textbf{TRENNBARE VERBEN}}\\
& \textbf{---} & \\
\end{tabular}
\vspace{-5pt}
\end{verbentry}

% schwitzen (to sweat)
\begin{verbentry}{
  style = {default},
  toc = {schwitzen (to sweat)},
  name = {schwitzen},
  english = {to sweat},
  type = {regular},
  auxiliary = {haben}
}


\vspace{5pt}
\hrule
\vspace{5pt}

\begin{tabular}{rl}
\multicolumn{2}{l}{\textbf{PRÄSENS}}\\
ich & \textbf{schwitze}\\
du & \textbf{schwitzt}\\
er/sie/es & \textbf{schwitzt}\\
wir & \textbf{schwitzen}\\
ihr & \textbf{schwitzt}\\
sie/Sie & \textbf{schwitzen}\\
\end{tabular}
\hfill
\begin{tabular}{rl}
\multicolumn{2}{l}{\textbf{PRÄTERITUM}}\\
ich & \textbf{schwitzte}\\
du & \textbf{schwitztest}\\
er/sie/es & \textbf{schwitzte}\\
wir & \textbf{schwitzten}\\
ihr & \textbf{schwitztet}\\
sie/Sie & \textbf{schwitzten}\\
\end{tabular}
\hfill
\begin{tabular}{rl}
\multicolumn{2}{l}{\textbf{IMPERATIV}}\\
& \textbf{schwitze (du)!}\\
& \textbf{schwitzen wir!}\\
& \textbf{schwitzt (ihr)!}\\
& \textbf{schwitzen Sie!}\\
\end{tabular}

\vspace{10pt}

\begin{tabular}{rl}
\multicolumn{2}{l}{\textbf{PERFEKT}}\\
& haben + \textbf{geschwitzt}\\
\end{tabular}
\hfill
\begin{tabular}{rl}
\multicolumn{2}{l}{\textbf{FUTURE I}}\\
& werden + \textbf{schwitzen}\\
\end{tabular}
\hfill
\begin{tabular}{rl}
\multicolumn{2}{l}{\textbf{PLUSQUAMPERFEKT}}\\
& hatten + \textbf{geschwitzt}\\
\end{tabular}

\vspace{10pt}

\begin{tabular}{lll}
\multicolumn{3}{l}{\textbf{MIT PRÄPOSITIONEN}}\\
& \textbf{schwitzen bei} + Dat & (to sweat during)\\
& \textbf{schwitzen vor} + Dat & (to sweat from)\\
\end{tabular}
\hfill
\begin{tabular}{lll}
\multicolumn{3}{l}{\textbf{TRENNBARE VERBEN}}\\
& \textbf{---} & \\
\end{tabular}

\vspace{-5pt}
\end{verbentry}

% sehen (to see)
\begin{verbentry}{
  style = {acc},
  toc = {sehen (to see)},
  name = {sehen},
  english = {to see},
  type = {irregular, + Akkusativ},
  auxiliary = {haben}
}
 
    
     \vspace{5pt}

    \hrule
    
    \vspace{5pt}

  \begin{tabular}{rl}
     \multicolumn{2}{l}{\textbf{PRÄSENS} }\\
    ich & \textbf{sehe}\\
    du & \textbf{siehst}\\
    er/sie/es & \textbf{sieht}\\
    wir & \textbf{sehen}\\
    ihr & \textbf{seht}\\
    sie/Sie & \textbf{sehen}\\
  \end{tabular}
  \hfill
     \begin{tabular}{rl}
    \multicolumn{2}{l}{\textbf{PRÄTERITUM} }\\
    ich & \textbf{sah}\\
    du & \textbf{sahst}\\
    er/sie/es & \textbf{sah}\\
    wir & \textbf{sahen}\\
    ihr & \textbf{saht}\\
    sie/Sie & \textbf{sahen}\\
  \end{tabular}
  \hfill
  \begin{tabular}{rl}
     \multicolumn{2}{l}{\textbf{IMPERATIV} }\\
   & \textbf{sieh(e) (du)!}\\
   & \textbf{sehen wir!}\\
   & \textbf{seht (ihr)!}\\
   & \textbf{sehen Sie!}\\
   & \vspace{10pt}
  \end{tabular}

    \vspace{10pt}

 \begin{tabular}{rl}
     \multicolumn{2}{l}{\textbf{PERFEKT}}\\
  &  haben + \textbf{gesehen}\\
  \end{tabular}
\hfill
   \begin{tabular}{rl}
   
    \multicolumn{2}{l}{\textbf{FUTURE I} }\\
  &  werden + \textbf{sehen}\\
 
  \end{tabular}
\hfill
   \begin{tabular}{rl}
      \multicolumn{2}{l}{\textbf{PLUSQUAMPERFEKT}}\\
  &  hatten + \textbf{gesehen}\\
  \end{tabular}
  


 \vspace{10pt}

   \begin{tabular}{lll}
      \multicolumn{3}{l}{\textbf{MIT PRÄPOSITIONEN}}\\
& \textbf{sehen auf} + Akk & (to pay attention to)\\
& \textbf{sehen in} + Akk & (to look into)\\
& \textbf{sehen nach} + Dat & (to check / look after)\\
& \textbf{absehen von} + Dat & (to refrain from)\\
& \textbf{ansehen als} + Akk & (to regard as)\\
& \textbf{ansehen von} + Dat & (to look from)\\
& \textbf{hinsehen auf} + Akk & (to look at)\\
& \textbf{wegsehen bei} + Dat & (to ignore during)\\
& \textbf{zusehen bei} + Dat & (to watch while)\\
  \end{tabular}
  \hfill
   \begin{tabular}{lll}
      \multicolumn{3}{l}{\textbf{TRENNBARE VERBEN}}\\
 & \textbf{an$|$sehen} & (to regard as)\\
  & \textbf{ab$|$sehen} & (to foresee / refrain)\\
   & \textbf{auf$|$sehen} & (to look up)\\
     & \textbf{aus$|$sehen} & (to appear)\\
    & \textbf{durch$|$sehen} & (to look through / examine)\\
      & \textbf{ein$|$sehen} & (to review)\\
    & \textbf{hin$|$sehen} & (to look (in a direction))\\
      & \textbf{um$|$sehen} & (to look around)\\
     & \textbf{weg$|$sehen} & (to look away / ignore)\\
      & \textbf{zu$|$sehen} & (to watch)\\
  \end{tabular}


  \vspace{-5pt}
\end{verbentry}

% setzen (to set / place / put)
\begin{verbentry}{
  style = {default},
  toc = {setzen (to set / place / put)},
  name = {setzen},
  english = {to set / place / put},
  type = {regular},
  auxiliary = {haben}
}
 
    
     \vspace{5pt}

    \hrule
    
    \vspace{5pt}

  \begin{tabular}{rl}
     \multicolumn{2}{l}{\textbf{PRÄSENS} }\\
    ich & \textbf{setze}\\
    du & \textbf{setzt}\\
    er/sie/es & \textbf{setzt}\\
    wir & \textbf{setzen}\\
    ihr & \textbf{setzt}\\
    sie/Sie & \textbf{setzen}\\
  \end{tabular}
  \hfill
  \begin{tabular}{rl}
    \multicolumn{2}{l}{\textbf{PRÄTERITUM} }\\
    ich & \textbf{setzte}\\
    du & \textbf{setztest}\\
    er/sie/es & \textbf{setzte}\\
    wir & \textbf{setzten}\\
    ihr & \textbf{setztet}\\
    sie/Sie & \textbf{setzten}\\
  \end{tabular}
  \hfill
  \begin{tabular}{rl}
     \multicolumn{2}{l}{\textbf{IMPERATIV} }\\
   & \textbf{setze (du)!}\\
   & \textbf{setzen wir!}\\
   & \textbf{setzt (ihr)!}\\
   & \textbf{setzen Sie!}\\
   & \vspace{10pt}
  \end{tabular}

 \vspace{10pt}

 \begin{tabular}{rl}
     \multicolumn{2}{l}{\textbf{PERFEKT}}\\
   & haben + \textbf{gesetzt}\\
  \end{tabular}
\hfill
\begin{tabular}{rl}
    \multicolumn{2}{l}{\textbf{FUTURE I}}\\
   & werden + \textbf{setzen}\\
  \end{tabular}
\hfill
\begin{tabular}{rl}
    \multicolumn{2}{l}{\textbf{PLUSQUAMPERFEKT}}\\
   & hatten + \textbf{gesetzt}\\
  \end{tabular}

 \vspace{10pt}

\begin{tabular}{lll}
      \multicolumn{3}{l}{\textbf{MIT PRÄPOSITIONEN}}\\
& \textbf{setzen auf} + Akk & (to place on / set on)\\
& \textbf{setzen in} + Akk & (to place into / put in)\\
& \textbf{sich setzen auf} + Akk & (to sit down on)\\
& \textbf{setzen für} + Akk & (to set for / dedicate)\\
\end{tabular}
  \hfill
\begin{tabular}{lll}
      \multicolumn{3}{l}{\textbf{TRENNBARE VERBEN}}\\
 & \textbf{ein$|$setzen} & (to insert / deploy / put in)\\
 & \textbf{auf$|$setzen} & (to place on / put onto)\\
 & \textbf{durch$|$setzen} & (to enforce / push through)\\
\vspace{10pt}
  \end{tabular}


\vspace{-5pt}
\end{verbentry}

% siezen (to address someone with ``Sie")
\begin{verbentry}{
  style = {default},
  toc = {siezen (to address someone with ``Sie")},
  name = {siezen},
  english = {to address someone with \enquote{Sie}},
  type = {regular},
  auxiliary = {haben}
}


\vspace{5pt}
\hrule
\vspace{5pt}

\begin{tabular}{rl}
\multicolumn{2}{l}{\textbf{PRÄSENS}}\\
ich & \textbf{sieze}\\
du & \textbf{siezt}\\
er/sie/es & \textbf{siezt}\\
wir & \textbf{siezen}\\
ihr & \textbf{siezt}\\
sie/Sie & \textbf{siezen}\\
\end{tabular}
\hfill
\begin{tabular}{rl}
\multicolumn{2}{l}{\textbf{PRÄTERITUM}}\\
ich & \textbf{siezte}\\
du & \textbf{sieztest}\\
er/sie/es & \textbf{siezte}\\
wir & \textbf{siezten}\\
ihr & \textbf{sieztet}\\
sie/Sie & \textbf{siezten}\\
\end{tabular}
\hfill
\begin{tabular}{rl}
\multicolumn{2}{l}{\textbf{IMPERATIV}}\\
& \textbf{sieze (du)!}\\
& \textbf{siezen wir!}\\
& \textbf{siezt (ihr)!}\\
& \textbf{siezen Sie!}\\
\end{tabular}

\vspace{10pt}

\begin{tabular}{rl}
\multicolumn{2}{l}{\textbf{PERFEKT}}\\
& haben + \textbf{gesiezt}\\
\end{tabular}
\hfill
\begin{tabular}{rl}
\multicolumn{2}{l}{\textbf{FUTURE I}}\\
& werden + \textbf{siezen}\\
\end{tabular}
\hfill
\begin{tabular}{rl}
\multicolumn{2}{l}{\textbf{PLUSQUAMPERFEKT}}\\
& hatten + \textbf{gesiezt}\\
\end{tabular}

\vspace{10pt}

\begin{tabular}{lll}
\multicolumn{3}{l}{\textbf{MIT PRÄPOSITIONEN}}\\
& \textbf{siezen mit} + Dat & (to use "Sie" with someone)\\
\end{tabular}
\hfill
\begin{tabular}{lll}
\multicolumn{3}{l}{\textbf{TRENNBARE VERBEN}}\\
& \textbf{---} & \\
\end{tabular}
\vspace{-5pt}
\end{verbentry}

% sitzen (to sit)
\begin{verbentry}{
  style = {default},
  toc = {sitzen (to sit)},
  name = {sitzen},
  english = {to sit},
  type = {irregular},
  auxiliary = {haben / sein context-dependent}
}
 
    
     \vspace{5pt}

    \hrule
    
    \vspace{5pt}

  \begin{tabular}{rl}
     \multicolumn{2}{l}{\textbf{PRÄSENS} }\\
    ich & \textbf{sitze}\\
    du & \textbf{sitzt}\\
    er/sie/es & \textbf{sitzt}\\
    wir & \textbf{sitzen}\\
    ihr & \textbf{sitzt}\\
    sie/Sie & \textbf{sitzen}\\
  \end{tabular}
  \hfill
  \begin{tabular}{rl}
    \multicolumn{2}{l}{\textbf{PRÄTERITUM} }\\
    ich & \textbf{saß}\\
    du & \textbf{saßt}\\
    er/sie/es & \textbf{saß}\\
    wir & \textbf{saßen}\\
    ihr & \textbf{saßt}\\
    sie/Sie & \textbf{saßen}\\
  \end{tabular}
  \hfill
  \begin{tabular}{rl}
     \multicolumn{2}{l}{\textbf{IMPERATIV} }\\
   & \textbf{sitze (du)!}\\
   & \textbf{sitzen wir!}\\
   & \textbf{sitzt (ihr)!}\\
   & \textbf{sitzen Sie!}\\
   & \vspace{10pt}
  \end{tabular}

 \vspace{10pt}

 \begin{tabular}{rl}
     \multicolumn{2}{l}{\textbf{PERFEKT}}\\
   & haben + \textbf{gesessen}\\
  \end{tabular}
\hfill
\begin{tabular}{rl}
    \multicolumn{2}{l}{\textbf{FUTURE I}}\\
   & werden + \textbf{sitzen}\\
  \end{tabular}
\hfill
\begin{tabular}{rl}
    \multicolumn{2}{l}{\textbf{PLUSQUAMPERFEKT}}\\
   & hatten + \textbf{gesessen}\\
  \end{tabular}

 \vspace{10pt}

\begin{tabular}{lll}
      \multicolumn{3}{l}{\textbf{MIT PRÄPOSITIONEN}}\\
& \textbf{sitzen auf} + Dat & (to sit on)\\
& \textbf{sitzen in} + Dat & (to sit in / inside)\\
& \textbf{sitzen an} + Dat & (to sit at / by)\\
\end{tabular}
  \hfill
\begin{tabular}{lll}
      \multicolumn{3}{l}{\textbf{TRENNBARE VERBEN}}\\
 & \textbf{abs$|$sitzen} & (to be seated / sit down)\\
 & \textbf{zusammen$|$sitzen} & (to sit together / be grouped)\\
 & \textbf{nach$|$sitzen} & (to stay after school / serve detention)\\
  \end{tabular}
  \vspace{-5pt}
\end{verbentry}

% sparen (to save)
\begin{verbentry}{
  style = {acc},
  toc = {sparen (to save)},
  name = {sparen},
  english = {to save},
  type = {regular, + Akkusativ},
  auxiliary = {haben}
}
 
    
     \vspace{5pt}

    \hrule
    
    \vspace{5pt}

  \begin{tabular}{rl}
     \multicolumn{2}{l}{\textbf{PRÄSENS} }\\
    ich & \textbf{spare}\\
    du & \textbf{sparst}\\
    er/sie/es & \textbf{spart}\\
    wir & \textbf{sparen}\\
    ihr & \textbf{spart}\\
    sie/Sie & \textbf{sparen}\\
  \end{tabular}
  \hfill
     \begin{tabular}{rl}
    \multicolumn{2}{l}{\textbf{PRÄTERITUM} }\\
    ich & \textbf{sparte}\\
    du & \textbf{spartest}\\
    er/sie/es & \textbf{sparte}\\
    wir & \textbf{sparten}\\
    ihr & \textbf{spartet}\\
    sie/Sie & \textbf{sparten}\\
  \end{tabular}
  \hfill
  \begin{tabular}{rl}
     \multicolumn{2}{l}{\textbf{IMPERATIV} }\\
   & \textbf{spar(e) (du)!}\\
   & \textbf{sparen wir!}\\
   & \textbf{sparet (ihr)!}\\
   & \textbf{sparen Sie!}\\
   & \vspace{10pt}
  \end{tabular}

    \vspace{10pt}

 \begin{tabular}{rl}
     \multicolumn{2}{l}{\textbf{PERFEKT}}\\
  &  haben + \textbf{gespart}\\
  \end{tabular}
\hfill
   \begin{tabular}{rl}
   
    \multicolumn{2}{l}{\textbf{FUTURE I} }\\
  &  werden + \textbf{sparen}\\
 
  \end{tabular}
\hfill
   \begin{tabular}{rl}
      \multicolumn{2}{l}{\textbf{PLUSQUAMPERFEKT}}\\
  &  hatten + \textbf{gesparet}\\
  \end{tabular}
  
  
\vspace{10pt}

\begin{tabular}{lll}
\multicolumn{3}{l}{\textbf{MIT PRÄPOSITIONEN}}\\
& \textbf{---} & \\
\end{tabular}
\hfill
\begin{tabular}{lll}
\multicolumn{3}{l}{\textbf{TRENNBARE VERBEN}}\\
& \textbf{---} & \\
\end{tabular}
\vspace{-5pt}
\end{verbentry}

% spielen (to play)
\begin{verbentry}{
  style = {acc},
  toc = {spielen (to play)},
  name = {spielen},
  english = {to play},
  type = {regular, + Akkusativ},
  auxiliary = {haben}
}
 
    
     \vspace{5pt}

    \hrule
    
    \vspace{5pt}

  \begin{tabular}{rl}
     \multicolumn{2}{l}{\textbf{PRÄSENS} }\\
    ich & \textbf{spiele}\\
    du & \textbf{spielst}\\
    er/sie/es & \textbf{spielt}\\
    wir & \textbf{spielen}\\
    ihr & \textbf{spielt}\\
    sie/Sie & \textbf{spielen}\\
  \end{tabular}
  \hfill
     \begin{tabular}{rl}
    \multicolumn{2}{l}{\textbf{PRÄTERITUM} }\\
    ich & \textbf{spielte}\\
    du & \textbf{spieltest}\\
    er/sie/es & \textbf{spielte}\\
    wir & \textbf{spielten}\\
    ihr & \textbf{spieltet}\\
    sie/Sie & \textbf{spielten}\\
  \end{tabular}
  \hfill
  \begin{tabular}{rl}
     \multicolumn{2}{l}{\textbf{IMPERATIV} }\\
   & \textbf{spiel(e) (du)!}\\
   & \textbf{spielen wir!}\\
   & \textbf{spielt (ihr)!}\\
   & \textbf{spielen Sie!}\\
   & \vspace{10pt}
  \end{tabular}

    \vspace{10pt}

 \begin{tabular}{rl}
     \multicolumn{2}{l}{\textbf{PERFEKT}}\\
   & haben + \textbf{gespielt}\\
  \end{tabular}
\hfill
   \begin{tabular}{rl}
   
    \multicolumn{2}{l}{\textbf{FUTURE I} }\\
  &  werden + \textbf{spielen}\\
 
  \end{tabular}
\hfill
   \begin{tabular}{rl}
      \multicolumn{2}{l}{\textbf{PLUSQUAMPERFEKT}}\\
   & hatten + \textbf{gespielt}\\
  \end{tabular}
  
  
\vspace{10pt}

\begin{tabular}{lll}
\multicolumn{3}{l}{\textbf{MIT PRÄPOSITIONEN}}\\
& \textbf{---} & \\
\end{tabular}
\hfill
\begin{tabular}{lll}
\multicolumn{3}{l}{\textbf{TRENNBARE VERBEN}}\\
& \textbf{---} & \\
\end{tabular}
\vspace{-5pt}
\end{verbentry}

% sprechen (to speak)
\begin{verbentry}{
  style = {default},
  toc = {sprechen (to speak)},
  name = {sprechen},
  english = {to speak},
  type = {irregular},
  auxiliary = {haben}
}
 
    
     \vspace{5pt}

    \hrule
    
    \vspace{5pt}

  \begin{tabular}{rl}
     \multicolumn{2}{l}{\textbf{PRÄSENS} }\\
    ich & \textbf{spreche}\\
    du & \textbf{sprichst}\\
    er/sie/es & \textbf{spricht}\\
    wir & \textbf{sprechen}\\
    ihr & \textbf{sprecht}\\
    sie/Sie & \textbf{sprechen}\\
  \end{tabular}
  \hfill
     \begin{tabular}{rl}
    \multicolumn{2}{l}{\textbf{PRÄTERITUM} }\\
    ich & \textbf{sprach}\\
    du & \textbf{sprachst}\\
    er/sie/es & \textbf{sprach}\\
    wir & \textbf{sprachen}\\
    ihr & \textbf{spracht}\\
    sie/Sie & \textbf{sprachen}\\
  \end{tabular}
  \hfill
  \begin{tabular}{rl}
     \multicolumn{2}{l}{\textbf{IMPERATIV} }\\
   & \textbf{sprich (du)!}\\
   & \textbf{sprechen wir!}\\
   & \textbf{sprecht (ihr)!}\\
   & \textbf{sprechen Sie!}\\
   & \vspace{10pt}
  \end{tabular}

    \vspace{10pt}

 \begin{tabular}{rl}
     \multicolumn{2}{l}{\textbf{PERFEKT}}\\
  &  haben + \textbf{gesprochen}\\
  \end{tabular}
\hfill
   \begin{tabular}{rl}
   
    \multicolumn{2}{l}{\textbf{FUTURE I} }\\
   & werden + \textbf{sprechen}\\
 
  \end{tabular}
\hfill
   \begin{tabular}{rl}
      \multicolumn{2}{l}{\textbf{PLUSQUAMPERFEKT}}\\
   & hatten + \textbf{gesprochen}\\
  \end{tabular}
  
   \vspace{10pt}

   \begin{tabular}{lll}
      \multicolumn{3}{l}{\textbf{MIT PRÄPOSITIONEN}}\\
&\textbf{sprechen für} + Akk & (to speak for)\\
& \textbf{sprechen über} + Akk & (to speak about)\\
& \textbf{sprechen mit} + Dat & (to speak with)\\
&\textbf{sprechen von} + Dat & (to speak of)
  \end{tabular}
  \hfill
   \begin{tabular}{lll}
      \multicolumn{3}{l}{\textbf{TRENNBARE VERBEN}}\\
 & \textbf{aus$|$sprechen} & (to pronounce)\\
&\textbf{ab$|$sprechen} & (to agree / cancel)\\
&\textbf{nach$|$sprechen} & (to repeat)
\vspace{10pt}
  \end{tabular}
  

  \vspace{-5pt}
\end{verbentry}

% stapeln (to stack / pile up)
\begin{verbentry}{
  style = {default},
  toc = {stapeln (to stack / pile up)},
  name = {stapeln},
  english = {to stack / pile up},
  type = {regular},
  auxiliary = {haben}
}


\vspace{5pt}
\hrule
\vspace{5pt}

\begin{tabular}{rl}
\multicolumn{2}{l}{\textbf{PRÄSENS}}\\
ich & \textbf{staple}\\
du & \textbf{stapelst}\\
er/sie/es & \textbf{stapelt}\\
wir & \textbf{stapeln}\\
ihr & \textbf{stapelt}\\
sie/Sie & \textbf{stapeln}\\
\end{tabular}
\hfill
\begin{tabular}{rl}
\multicolumn{2}{l}{\textbf{PRÄTERITUM}}\\
ich & \textbf{stapelte}\\
du & \textbf{stapeltest}\\
er/sie/es & \textbf{stapelte}\\
wir & \textbf{stapelten}\\
ihr & \textbf{stapeltet}\\
sie/Sie & \textbf{stapelten}\\
\end{tabular}
\hfill
\begin{tabular}{rl}
\multicolumn{2}{l}{\textbf{IMPERATIV}}\\
& \textbf{staple (du)!}\\
& \textbf{stapeln wir!}\\
& \textbf{stapelt (ihr)!}\\
& \textbf{stapeln Sie!}\\
\end{tabular}

\vspace{10pt}

\begin{tabular}{rl}
\multicolumn{2}{l}{\textbf{PERFEKT}}\\
& haben + \textbf{gestapelt}\\
\end{tabular}
\hfill
\begin{tabular}{rl}
\multicolumn{2}{l}{\textbf{FUTURE I}}\\
& werden + \textbf{stapeln}\\
\end{tabular}
\hfill
\begin{tabular}{rl}
\multicolumn{2}{l}{\textbf{PLUSQUAMPERFEKT}}\\
& hatten + \textbf{gestapelt}\\
\end{tabular}

\vspace{10pt}

\begin{tabular}{lll}
\multicolumn{3}{l}{\textbf{MIT PRÄPOSITIONEN}}\\
& \textbf{stapeln auf} + Akk & (to stack onto)\\
& \textbf{stapeln in} + Akk & (to stack into)\\
\end{tabular}
\hfill
\begin{tabular}{lll}
\multicolumn{3}{l}{\textbf{TRENNBARE VERBEN}}\\
& \textbf{auf$|$stapeln} & (to stack up)\\
\end{tabular}
\vspace{-5pt}
\end{verbentry}

% starten (to start / launch)
\begin{verbentry}{
  style = {default},
  toc = {starten (to start / launch)},
  name = {starten},
  english = {to start / launch},
  type = {regular},
  auxiliary = {haben}
}
 
    
     \vspace{5pt}

    \hrule
    
    \vspace{5pt}

  \begin{tabular}{rl}
     \multicolumn{2}{l}{\textbf{PRÄSENS} }\\
    ich & \textbf{starte}\\
    du & \textbf{startest}\\
    er/sie/es & \textbf{startet}\\
    wir & \textbf{starten}\\
    ihr & \textbf{startet}\\
    sie/Sie & \textbf{starten}\\
  \end{tabular}
  \hfill
  \begin{tabular}{rl}
    \multicolumn{2}{l}{\textbf{PRÄTERITUM} }\\
    ich & \textbf{startete}\\
    du & \textbf{startetest}\\
    er/sie/es & \textbf{startete}\\
    wir & \textbf{starteten}\\
    ihr & \textbf{startetet}\\
    sie/Sie & \textbf{starteten}\\
  \end{tabular}
  \hfill
  \begin{tabular}{rl}
     \multicolumn{2}{l}{\textbf{IMPERATIV} }\\
   & \textbf{starte (du)!}\\
   & \textbf{starten wir!}\\
   & \textbf{startet (ihr)!}\\
   & \textbf{starten Sie!}\\
   & \vspace{10pt}
  \end{tabular}

 \vspace{10pt}

 \begin{tabular}{rl}
     \multicolumn{2}{l}{\textbf{PERFEKT}}\\
   & haben + \textbf{gestartet}\\
  \end{tabular}
\hfill
\begin{tabular}{rl}
    \multicolumn{2}{l}{\textbf{FUTURE I}}\\
   & werden + \textbf{starten}\\
  \end{tabular}
\hfill
\begin{tabular}{rl}
    \multicolumn{2}{l}{\textbf{PLUSQUAMPERFEKT}}\\
   & hatten + \textbf{gestartet}\\
  \end{tabular}

 \vspace{10pt}

\begin{tabular}{lll}
      \multicolumn{3}{l}{\textbf{MIT PRÄPOSITIONEN}}\\
& \textbf{starten zu} + Dat & (to start to / head towards)\\
& \textbf{starten mit} + Dat & (to start with)\\
& \textbf{starten von} + Dat & (to start from)\\
\end{tabular}
  \hfill
\begin{tabular}{lll}
      \multicolumn{3}{l}{\textbf{TRENNBARE VERBEN}}\\
 & \textbf{ab$|$starten} & (to take off / depart)\\
 & \textbf{auf$|$starten} & (to launch / initiate)\\
\vspace{10pt}
  \end{tabular}


\vspace{-5pt}
\end{verbentry}

% staubsaugen (to vacuum)
\begin{verbentry}{
  style = {default},
  toc = {staubsaugen (to vacuum)},
  name = {staubsaugen},
  english = {to vacuum},
  type = {mixed / separable},
  auxiliary = {haben}
}
 

\vspace{5pt}
\hrule
\vspace{5pt}

\begin{tabular}{rl}
\multicolumn{2}{l}{\textbf{PRÄSENS}}\\
ich & \textbf{sauge Staub}\\
du & \textbf{saugst Staub}\\
er/sie/es & \textbf{saugt Staub}\\
wir & \textbf{saugen Staub}\\
ihr & \textbf{saugt Staub}\\
sie/Sie & \textbf{saugen Staub}\\
\end{tabular}
\hfill
\begin{tabular}{rl}
\multicolumn{2}{l}{\textbf{PRÄTERITUM}}\\
ich & \textbf{saugte Staub}\\
du & \textbf{saugtest Staub}\\
er/sie/es & \textbf{saugte Staub}\\
wir & \textbf{saugten Staub}\\
ihr & \textbf{saugtet Staub}\\
sie/Sie & \textbf{saugten Staub}\\
\end{tabular}
\hfill
\begin{tabular}{rl}
\multicolumn{2}{l}{\textbf{IMPERATIV}}\\
& \textbf{saug Staub (du)!}\\
& \textbf{saugen wir Staub!}\\
& \textbf{saugt Staub (ihr)!}\\
& \textbf{saugen Sie Staub!}\\
\end{tabular}

\vspace{10pt}

\begin{tabular}{rl}
\multicolumn{2}{l}{\textbf{PERFEKT}}\\
& haben + \textbf{staubgesaugt}\\
\end{tabular}
\hfill
\begin{tabular}{rl}
\multicolumn{2}{l}{\textbf{FUTURE I}}\\
& werden + \textbf{staubsaugen}\\
\end{tabular}
\hfill
\begin{tabular}{rl}
\multicolumn{2}{l}{\textbf{PLUSQUAMPERFEKT}}\\
& hatten + \textbf{staubgesaugt}\\
\end{tabular}

\vspace{10pt}

\begin{tabular}{lll}
\multicolumn{3}{l}{\textbf{MIT PRÄPOSITIONEN}}\\
& \textbf{staubsaugen} + Akk & (to vacuum something)\\
\end{tabular}
\hfill
\begin{tabular}{lll}
\multicolumn{3}{l}{\textbf{TRENNBARE VERBEN}}\\
& \textbf{---} & \\
\end{tabular}

\vspace{-5pt}
\end{verbentry}

% stechen (to sting / to stab)
\begin{verbentry}{
  style = {default},
  toc = {stechen (to sting / to stab)},
  name = {stechen},
  english = {to sting / to stab},
  type = {irregular},
  auxiliary = {haben}
}
 
    
     \vspace{5pt}

    \hrule
    
    \vspace{5pt}

  \begin{tabular}{rl}
     \multicolumn{2}{l}{\textbf{PRÄSENS} }\\
    ich & \textbf{steche}\\
    du & \textbf{stichst}\\
    er/sie/es & \textbf{sticht}\\
    wir & \textbf{stechen}\\
    ihr & \textbf{stecht}\\
    sie/Sie & \textbf{stechen}\\
  \end{tabular}
  \hfill
     \begin{tabular}{rl}
    \multicolumn{2}{l}{\textbf{PRÄTERITUM} }\\
    ich & \textbf{stach}\\
    du & \textbf{stachst}\\
    er/sie/es & \textbf{stach}\\
    wir & \textbf{stachen}\\
    ihr & \textbf{stacht}\\
    sie/Sie & \textbf{stachen}\\
  \end{tabular}
  \hfill
  \begin{tabular}{rl}
     \multicolumn{2}{l}{\textbf{IMPERATIV} }\\
   & \textbf{stich (du)!}\\
   & \textbf{stechen wir!}\\
   & \textbf{stecht (ihr)!}\\
   & \textbf{stechen Sie!}\\
   & \vspace{10pt}
  \end{tabular}

    \vspace{10pt}

 \begin{tabular}{rl}
     \multicolumn{2}{l}{\textbf{PERFEKT}}\\
   & haben + \textbf{gestochen}\\
  \end{tabular}
\hfill
   \begin{tabular}{rl}
    \multicolumn{2}{l}{\textbf{FUTURE I} }\\
   & werde + \textbf{stechen}\\
  \end{tabular}
\hfill
   \begin{tabular}{rl}
      \multicolumn{2}{l}{\textbf{PLUSQUAMPERFEKT}}\\
   & hatten + \textbf{gestochen}\\
  \end{tabular}

   \vspace{10pt}

   \begin{tabular}{lll}
      \multicolumn{3}{l}{\textbf{MIT PRÄPOSITIONEN}}\\
& \textbf{stechen in} + Akk & (to sting/stab into)\\
& \textbf{stechen mit} + Dat & (to sting/stab with)\\
\end{tabular}
  \hfill
   \begin{tabular}{lll}
      \multicolumn{3}{l}{\textbf{TRENNBARE VERBEN}}\\
 & \textbf{zu$|$stechen} & (to stab)\\
 & \textbf{ein$|$stechen} & (to stab into / pierce)\\
\vspace{10pt}
  \end{tabular}

  \vspace{-5pt}
\end{verbentry}

% stehen (to stand)
\begin{verbentry}{
  style = {default},
  toc = {stehen (to stand)},
  name = {stehen},
  english = {to stand},
  type = {irregular},
  auxiliary = {haben/sein}
}
 
    
     \vspace{5pt}

    \hrule
    
    \vspace{5pt}

  \begin{tabular}{rl}
     \multicolumn{2}{l}{\textbf{PRÄSENS} }\\
    ich & \textbf{stehe}\\
    du & \textbf{stehst}\\
    er/sie/es & \textbf{steht}\\
    wir & \textbf{stehen}\\
    ihr & \textbf{steht}\\
    sie/Sie & \textbf{stehen}\\
  \end{tabular}
  \hfill
     \begin{tabular}{rl}
    \multicolumn{2}{l}{\textbf{PRÄTERITUM} }\\
    ich & \textbf{stand}\\
    du & \textbf{standst}\\
    er/sie/es & \textbf{stand}\\
    wir & \textbf{standen}\\
    ihr & \textbf{standet}\\
    sie/Sie & \textbf{standen}\\
  \end{tabular}
  \hfill
  \begin{tabular}{rl}
     \multicolumn{2}{l}{\textbf{IMPERATIV} }\\
   & \textbf{steh(e) (du)!}\\
   & \textbf{stehen wir!}\\
   & \textbf{steht (ihr)!}\\
   & \textbf{stehen Sie!}\\
   & \vspace{10pt}
  \end{tabular}

    \vspace{10pt}

 \begin{tabular}{rl}
     \multicolumn{2}{l}{\textbf{PERFEKT}}\\
  &  haben / sein + \textbf{gestanden}\\
  \end{tabular}
\hfill
   \begin{tabular}{rl}
   
    \multicolumn{2}{l}{\textbf{FUTURE I} }\\
  &  werden + \textbf{stehen}\\
 
  \end{tabular}
\hfill
   \begin{tabular}{rl}
      \multicolumn{2}{l}{\textbf{PLUSQUAMPERFEKT}}\\
  &  hatten / waren + \textbf{gestanden}\\
  \end{tabular}

   \vspace{10pt}

   \begin{tabular}{lll}
      \multicolumn{3}{l}{\textbf{MIT PRÄPOSITIONEN}}\\
& \textbf{stehen für} + Akk & (to represent)\\
& \textbf{stehen zu} + Dat & (to stand by)\\
&\textbf{stehen vor} + Dat & (to stand facing / in front of)\\
&\textbf{stehen in} + Dat & (to stand in)\\
&\textbf{stehen auf} + Dat & (to stand on)\\
&\textbf{stehen an} + Dat & (to stand at)\\
&\textbf{stehen neben} + Dat & (to stand next to)\\
&\textbf{stehen mit} + Dat & (to be connected with)\\
  \end{tabular}
  \hfill
   \begin{tabular}{lll}
      \multicolumn{3}{l}{\textbf{TRENNBARE VERBEN}}\\
 & \textbf{auf$|$stehen} & (to stand up, get up)\\
 & \textbf{an$|$stehen} & (to queue)\\
     & \textbf{fest$|$stehen} & (to be certain)\\
     & \textbf{zu$|$stehen} & (to be entitled to)\\
\vspace{40pt}
  \end{tabular}


  \vspace{-5pt}
\end{verbentry}

% steigen (to climb / rise)
\begin{verbentry}{
  style = {default},
  toc = {steigen (to climb / rise)},
  name = {steigen},
  english = {to climb / rise},
  type = {irregular},
  auxiliary = {sein}
}
 
    
     \vspace{5pt}

    \hrule
    
    \vspace{5pt}

  \begin{tabular}{rl}
     \multicolumn{2}{l}{\textbf{PRÄSENS} }\\
    ich & \textbf{steige}\\
    du & \textbf{steigst}\\
    er/sie/es & \textbf{steigt}\\
    wir & \textbf{steigen}\\
    ihr & \textbf{steigt}\\
    sie/Sie & \textbf{steigen}\\
  \end{tabular}
  \hfill
  \begin{tabular}{rl}
    \multicolumn{2}{l}{\textbf{PRÄTERITUM} }\\
    ich & \textbf{stieg}\\
    du & \textbf{stiegst}\\
    er/sie/es & \textbf{stieg}\\
    wir & \textbf{stiegen}\\
    ihr & \textbf{stiegt}\\
    sie/Sie & \textbf{stiegen}\\
  \end{tabular}
  \hfill
  \begin{tabular}{rl}
     \multicolumn{2}{l}{\textbf{IMPERATIV} }\\
   & \textbf{steige (du)!}\\
   & \textbf{steigen wir!}\\
   & \textbf{steigt (ihr)!}\\
   & \textbf{steigen Sie!}\\
   & \vspace{10pt}
  \end{tabular}

 \vspace{10pt}

 \begin{tabular}{rl}
     \multicolumn{2}{l}{\textbf{PERFEKT}}\\
   & sein + \textbf{gestiegen}\\
  \end{tabular}
\hfill
\begin{tabular}{rl}
    \multicolumn{2}{l}{\textbf{FUTURE I}}\\
   & werden + \textbf{steigen}\\
  \end{tabular}
\hfill
\begin{tabular}{rl}
    \multicolumn{2}{l}{\textbf{PLUSQUAMPERFEKT}}\\
   & waren + \textbf{gestiegen}\\
  \end{tabular}

 \vspace{10pt}

\begin{tabular}{lll}
      \multicolumn{3}{l}{\textbf{MIT PRÄPOSITIONEN}}\\
& \textbf{steigen auf} + Akk & (to climb onto / rise to)\\
& \textbf{steigen in} + Akk & (to climb into / rise into)\\
& \textbf{steigen aus} + Dat & (to climb out of)\\
\vspace{20pt}
\end{tabular}
  \hfill
\begin{tabular}{lll}
      \multicolumn{3}{l}{\textbf{TRENNBARE VERBEN}}\\
 & \textbf{auf$|$steigen} & (to take off / rise up)\\
 & \textbf{ein$|$steigen} & (to get on / board)\\
 & \textbf{aus$|$steigen} & (to get off / disembark)\\
 & \textbf{ab$|$steigen} & (to descend / go down)\\
  & \textbf{um$|$steigen} & (to transfer)\\
  \end{tabular}


  \vspace{-5pt}
\end{verbentry}

% stellen (to place)
\begin{verbentry}{
  style = {acc},
  toc = {stellen (to place)},
  name = {stellen},
  english = {to place},
  type = {regular, + Akkusativ},
  auxiliary = {haben}
}
 
    
     \vspace{5pt}

    \hrule
    
    \vspace{5pt}

  \begin{tabular}{rl}
     \multicolumn{2}{l}{\textbf{PRÄSENS} }\\
    ich & \textbf{stelle}\\
    du & \textbf{stellst}\\
    er/sie/es & \textbf{stellt}\\
    wir & \textbf{stellen}\\
    ihr & \textbf{stellt}\\
    sie/Sie & \textbf{stellen}\\
  \end{tabular}
  \hfill
     \begin{tabular}{rl}
    \multicolumn{2}{l}{\textbf{PRÄTERITUM} }\\
    ich & \textbf{stellte}\\
    du & \textbf{stelltest}\\
    er/sie/es & \textbf{stellte}\\
    wir & \textbf{stellten}\\
    ihr & \textbf{stelltet}\\
    sie/Sie & \textbf{stellten}\\
  \end{tabular}
  \hfill
  \begin{tabular}{rl}
     \multicolumn{2}{l}{\textbf{IMPERATIV} }\\
   & \textbf{stell(e) (du)!}\\
   & \textbf{stellen wir!}\\
   & \textbf{stellt (ihr)!}\\
   & \textbf{stellen Sie!}\\
   & \vspace{10pt}
  \end{tabular}

    \vspace{10pt}

 \begin{tabular}{rl}
     \multicolumn{2}{l}{\textbf{PERFEKT}}\\
   & haben + \textbf{gestellt}\\
  \end{tabular}
\hfill
   \begin{tabular}{rl}
   
    \multicolumn{2}{l}{\textbf{FUTURE I} }\\
   & werden + \textbf{stellen}\\
 
  \end{tabular}
\hfill
   \begin{tabular}{rl}
      \multicolumn{2}{l}{\textbf{PLUSQUAMPERFEKT}}\\
   & hatten + \textbf{gestellt}\\
  \end{tabular}
  
\vspace{10pt}

   \begin{tabular}{lll}
      \multicolumn{3}{l}{\textbf{MIT PRÄPOSITIONEN}}\\
&\textbf{sich stellen auf} + Akk & (to prepare for / brace for)\\
& \textbf{sich stellen gegen} + Akk & (to oppose)\\
& \textbf{sich stellen vor} + Akk & (to side with)\\
&\textbf{sich stellen zu} + Dat & (to stand in front of / defend)\\
\vspace{60pt}
  \end{tabular}
  \hfill
   \begin{tabular}{lll}
      \multicolumn{3}{l}{\textbf{TRENNBARE VERBEN}}\\
 & \textbf{auf$|$stellen} & (to set up / put up)\\
&\textbf{hin$|$stellen} & (to put (something) there)\\
&\textbf{um$|$stellen} & (to rearrange)\\
&\textbf{vor$|$stellen} & (to introduce)\\
&\textbf{ver$|$stellen} & (to block)\\
&\textbf{ein$|$stellen} & (to adjust, hire)\\
&\textbf{her$|$stellen} & (to produce)\\
&\textbf{dar$|$stellen} & (to portray)\\
&\textbf{fest$|$stellen} & (to determine, establish)\\
&\textbf{zurück$|$stellen} & (to postpone)
  \end{tabular}


  \vspace{-5pt}
\end{verbentry}

% stimmen (to be correct / to vote)
\begin{verbentry}{
  style = {dat},
  toc = {stimmen (to be correct / to vote)},
  name = {stimmen},
  english = {to be correct / to vote},
  type = {regular},
  auxiliary = {haben}
}
 
    
     \vspace{5pt}

    \hrule
    
    \vspace{5pt}

  \begin{tabular}{rl}
     \multicolumn{2}{l}{\textbf{PRÄSENS} }\\
    ich & \textbf{stimme}\\
    du & \textbf{stimmst}\\
    er/sie/es & \textbf{stimmt}\\
    wir & \textbf{stimmen}\\
    ihr & \textbf{stimmt}\\
    sie/Sie & \textbf{stimmen}\\
  \end{tabular}
  \hfill
     \begin{tabular}{rl}
    \multicolumn{2}{l}{\textbf{PRÄTERITUM} }\\
    ich & \textbf{stimmte}\\
    du & \textbf{stimmtest}\\
    er/sie/es & \textbf{stimmte}\\
    wir & \textbf{stimmten}\\
    ihr & \textbf{stimmtet}\\
    sie/Sie & \textbf{stimmten}\\
  \end{tabular}
  \hfill
  \begin{tabular}{rl}
     \multicolumn{2}{l}{\textbf{IMPERATIV} }\\
   & \textbf{stimm(e) (du)!}\\
   & \textbf{stimmen wir!}\\
   & \textbf{stimmt (ihr)!}\\
   & \textbf{stimmen Sie!}\\
   & \vspace{10pt}
  \end{tabular}

    \vspace{10pt}

 \begin{tabular}{rl}
     \multicolumn{2}{l}{\textbf{PERFEKT}}\\
   & haben + \textbf{gestimmt}\\
  \end{tabular}
\hfill
   \begin{tabular}{rl}
   
    \multicolumn{2}{l}{\textbf{FUTURE I} }\\
   & werden + \textbf{stimmen}\\
 
  \end{tabular}
\hfill
   \begin{tabular}{rl}
      \multicolumn{2}{l}{\textbf{PLUSQUAMPERFEKT}}\\
   & hatten + \textbf{gestimmt}\\
  \end{tabular}

   \vspace{10pt}

   \begin{tabular}{lll}
      \multicolumn{3}{l}{\textbf{MIT PRÄPOSITIONEN}}\\
& \textbf{stimmen für} + Akk & (to vote for)\\
& \textbf{stimmen gegen} + Akk & (to vote against)\\
\end{tabular}
  \hfill
   \begin{tabular}{lll}
      \multicolumn{3}{l}{\textbf{TRENNBARE VERBEN}}\\
 & \textbf{ab$|$stimmen} & (to vote / coordinate)\\
 & \textbf{überein$|$stimmen} & (to agree / correspond)\\
  \end{tabular}


  \vspace{-5pt}
\end{verbentry}

% stoppen (to stop)
\begin{verbentry}{
  style = {default},
  toc = {stoppen (to stop)},
  name = {stoppen},
  english = {to stop},
  type = {regular},
  auxiliary = {haben}
}
 
    
     \vspace{5pt}

    \hrule
    
    \vspace{5pt}

  \begin{tabular}{rl}
     \multicolumn{2}{l}{\textbf{PRÄSENS} }\\
    ich & \textbf{stoppe}\\
    du & \textbf{stoppst}\\
    er/sie/es & \textbf{stoppt}\\
    wir & \textbf{stoppen}\\
    ihr & \textbf{stoppt}\\
    sie/Sie & \textbf{stoppen}\\
  \end{tabular}
  \hfill
     \begin{tabular}{rl}
    \multicolumn{2}{l}{\textbf{PRÄTERITUM} }\\
    ich & \textbf{stoppte}\\
    du & \textbf{stopptest}\\
    er/sie/es & \textbf{stoppte}\\
    wir & \textbf{stoppten}\\
    ihr & \textbf{stopptet}\\
    sie/Sie & \textbf{stoppten}\\
  \end{tabular}
  \hfill
  \begin{tabular}{rl}
     \multicolumn{2}{l}{\textbf{IMPERATIV} }\\
   & \textbf{stopp(e) (du)!}\\
   & \textbf{stoppen wir!}\\
   & \textbf{stoppt (ihr)!}\\
   & \textbf{stoppen Sie!}\\
   & \vspace{10pt}
  \end{tabular}

    \vspace{10pt}

 \begin{tabular}{rl}
     \multicolumn{2}{l}{\textbf{PERFEKT}}\\
   & haben + \textbf{gestoppt}\\
  \end{tabular}
\hfill
   \begin{tabular}{rl}
    \multicolumn{2}{l}{\textbf{FUTURE I} }\\
   & werden + \textbf{stoppen}\\
  \end{tabular}
\hfill
   \begin{tabular}{rl}
      \multicolumn{2}{l}{\textbf{PLUSQUAMPERFEKT}}\\
   & hatten + \textbf{gestoppt}\\
  \end{tabular}

   \vspace{10pt}

   \begin{tabular}{lll}
      \multicolumn{3}{l}{\textbf{MIT PRÄPOSITIONEN}}\\
& \textbf{stoppen vor} + Dat & (to stop in front of)\\
& \textbf{stoppen bei} + Dat & (to stop at)\\
\end{tabular}
  \hfill
   \begin{tabular}{lll}
      \multicolumn{3}{l}{\textbf{TRENNBARE VERBEN}}\\
 & \textbf{an$|$stoppen} & (to stop / halt briefly)\\
\vspace{10pt}
  \end{tabular}

  \vspace{-5pt}
\end{verbentry}

% streichen (to paint / stroke)
\begin{verbentry}{
  style = {default},
  toc = {streichen (to paint / stroke)},
  name = {streichen},
  english = {to paint / stroke},
  type = {irregular},
  auxiliary = {haben}
}
 
    
     \vspace{5pt}

    \hrule
    
    \vspace{5pt}

  \begin{tabular}{rl}
     \multicolumn{2}{l}{\textbf{PRÄSENS} }\\
    ich & \textbf{streiche}\\
    du & \textbf{streichst}\\
    er/sie/es & \textbf{streicht}\\
    wir & \textbf{streichen}\\
    ihr & \textbf{streicht}\\
    sie/Sie & \textbf{streichen}\\
  \end{tabular}
  \hfill
  \begin{tabular}{rl}
    \multicolumn{2}{l}{\textbf{PRÄTERITUM} }\\
    ich & \textbf{strich}\\
    du & \textbf{strichst}\\
    er/sie/es & \textbf{strich}\\
    wir & \textbf{strichen}\\
    ihr & \textbf{stricht}\\
    sie/Sie & \textbf{strichen}\\
  \end{tabular}
  \hfill
  \begin{tabular}{rl}
     \multicolumn{2}{l}{\textbf{IMPERATIV} }\\
   & \textbf{streiche (du)!}\\
   & \textbf{streichen wir!}\\
   & \textbf{streicht (ihr)!}\\
   & \textbf{streichen Sie!}\\
   & \vspace{10pt}
  \end{tabular}

 \vspace{10pt}

 \begin{tabular}{rl}
     \multicolumn{2}{l}{\textbf{PERFEKT}}\\
   & haben + \textbf{gestrichen}\\
  \end{tabular}
\hfill
\begin{tabular}{rl}
    \multicolumn{2}{l}{\textbf{FUTURE I}}\\
   & werden + \textbf{streichen}\\
  \end{tabular}
\hfill
\begin{tabular}{rl}
    \multicolumn{2}{l}{\textbf{PLUSQUAMPERFEKT}}\\
   & hatten + \textbf{gestrichen}\\
  \end{tabular}

 \vspace{10pt}

\begin{tabular}{lll}
      \multicolumn{3}{l}{\textbf{MIT PRÄPOSITIONEN}}\\
& \textbf{streichen über} + Akk & (to stroke over / to paint over)\\
& \textbf{streichen mit} + Dat & (to paint / smear with)\\
\vspace{10pt}
\end{tabular}
  \hfill
\begin{tabular}{lll}
      \multicolumn{3}{l}{\textbf{TRENNBARE VERBEN}}\\
 & \textbf{ab$|$streichen} & (to cancel / remove / coat off)\\
 & \textbf{ein$|$streichen} & (to apply / cover in)\\
 & \textbf{auf$|$streichen} & (to spread / coat onto)\\
  \end{tabular}


\vspace{-5pt}
\end{verbentry}

% studieren (to study)
\begin{verbentry}{
  style = {default},
  toc = {studieren (to study)},
  name = {studieren},
  english = {to study},
  type = {regular},
  auxiliary = {haben}
}
 
    
     \vspace{5pt}

    \hrule
    
    \vspace{5pt}

  \begin{tabular}{rl}
     \multicolumn{2}{l}{\textbf{PRÄSENS} }\\
    ich & \textbf{studiere}\\
    du & \textbf{studierst}\\
    er/sie/es & \textbf{studiert}\\
    wir & \textbf{studieren}\\
    ihr & \textbf{studiert}\\
    sie/Sie & \textbf{studieren}\\
  \end{tabular}
  \hfill
     \begin{tabular}{rl}
    \multicolumn{2}{l}{\textbf{PRÄTERITUM} }\\
    ich & \textbf{studierte}\\
    du & \textbf{studiertest}\\
    er/sie/es & \textbf{studierte}\\
    wir & \textbf{studierten}\\
    ihr & \textbf{studiertet}\\
    sie/Sie & \textbf{studierten}\\
  \end{tabular}
  \hfill
  \begin{tabular}{rl}
     \multicolumn{2}{l}{\textbf{IMPERATIV} }\\
   & \textbf{studier(e) (du)!}\\
   & \textbf{studieren wir!}\\
   & \textbf{studiert (ihr)!}\\
   & \textbf{studieren Sie!}\\
   & \vspace{10pt}
  \end{tabular}

    \vspace{10pt}

 \begin{tabular}{rl}
     \multicolumn{2}{l}{\textbf{PERFEKT}}\\
   & haben + \textbf{studiert}\\
  \end{tabular}
\hfill
   \begin{tabular}{rl}
   
    \multicolumn{2}{l}{\textbf{FUTURE I} }\\
   & werden + \textbf{studieren}\\
 
  \end{tabular}
\hfill
   \begin{tabular}{rl}
      \multicolumn{2}{l}{\textbf{PLUSQUAMPERFEKT}}\\
   & hatten + \textbf{studiert}\\
  \end{tabular}
  
  \vspace{10pt}

\begin{tabular}{lll}
\multicolumn{3}{l}{\textbf{MIT PRÄPOSITIONEN}}\\
& \textbf{---} & \\
\end{tabular}
\hfill
\begin{tabular}{lll}
\multicolumn{3}{l}{\textbf{TRENNBARE VERBEN}}\\
& \textbf{---} & \\
\end{tabular}

  \vspace{-5pt}
\end{verbentry}

% suchen (to search)
\begin{verbentry}{
  style = {acc},
  toc = {suchen (to search)},
  name = {suchen},
  english = {to search},
  type = {regular, + Akkusativ},
  auxiliary = {haben}
}
 
    
     \vspace{5pt}

    \hrule
    
    \vspace{5pt}

  \begin{tabular}{rl}
     \multicolumn{2}{l}{\textbf{PRÄSENS} }\\
    ich & \textbf{suche}\\
    du & \textbf{suchst}\\
    er/sie/es & \textbf{sucht}\\
    wir & \textbf{suchen}\\
    ihr & \textbf{sucht}\\
    sie/Sie & \textbf{suchen}\\
  \end{tabular}
  \hfill
     \begin{tabular}{rl}
    \multicolumn{2}{l}{\textbf{PRÄTERITUM} }\\
    ich & \textbf{suchte}\\
    du & \textbf{suchtest}\\
    er/sie/es & \textbf{suchte}\\
    wir & \textbf{suchten}\\
    ihr & \textbf{suchtet}\\
    sie/Sie & \textbf{suchten}\\
  \end{tabular}
  \hfill
  \begin{tabular}{rl}
     \multicolumn{2}{l}{\textbf{IMPERATIV} }\\
   & \textbf{such(e) (du)!}\\
   & \textbf{suchen wir!}\\
   & \textbf{sucht (ihr)!}\\
   & \textbf{suchen Sie!}\\
   & \vspace{10pt}
  \end{tabular}

    \vspace{10pt}

 \begin{tabular}{rl}
     \multicolumn{2}{l}{\textbf{PERFEKT}}\\
   & haben + \textbf{gesucht}\\
  \end{tabular}
\hfill
   \begin{tabular}{rl}
   
    \multicolumn{2}{l}{\textbf{FUTURE I} }\\
   & werden + \textbf{suchen}\\
 
  \end{tabular}
\hfill
   \begin{tabular}{rl}
      \multicolumn{2}{l}{\textbf{PLUSQUAMPERFEKT}}\\
   & hatten + \textbf{gesucht}\\
  \end{tabular}
  
  
\vspace{10pt}

\begin{tabular}{lll}
\multicolumn{3}{l}{\textbf{MIT PRÄPOSITIONEN}}\\
& \textbf{---} & \\
\end{tabular}
\hfill
\begin{tabular}{lll}
\multicolumn{3}{l}{\textbf{TRENNBARE VERBEN}}\\
& \textbf{---} & \\
\end{tabular}
\vspace{-5pt}
\end{verbentry}
